\documentclass[12pt]{article}
\usepackage[T1]{fontenc}
\usepackage[utf8]{inputenc}
\usepackage{lmodern}
\usepackage{textcomp}
\usepackage{enumitem}
\usepackage{geometry}
\usepackage{graphicx}
\usepackage{amsmath, amssymb}
\geometry{margin=1in}
\begin{document}
\begin{center}
Engineering - Undergraduate Level Assessment 15 \
Total Marks: 40 \
Answer all questions.
\end{center}

\section*{Multiple Choice (10 marks)}
\begin{enumerate}[label=\arabic*.]
\item f(X) = [\ensuremath{\pi}(1 + $X^2$)]^-1- \ensuremath{\infty} \textless{} x \textless{} \ensuremath{\infty}. If Y = $X^2$, what is the density function of Y ?
\begin{enumerate}[label=(\alph*)]
    \item h(y) = [2 / \{\ensuremath{\pi}(1 + \ensuremath{\sqrt{y}})\}] for y \textgreater{} 0 and = 0 otherwise
    \item h(y) = [2 / \{\ensuremath{\pi}\ensuremath{\surdy}(1 + y)\}]y \textgreater{} 0 and = 0 otherwise
    \item h(y) = [1 / \{\ensuremath{\pi}(1 + $y^2$)\}] for y \textgreater{} 0 and = 0 otherwise
    \item h(y) = [1 / \{2\ensuremath{\pi}(1 + \ensuremath{\sqrt{y}})\}] for y \textgreater{} 0 and = 0 otherwise
    \item h(y) = [1 / \{2\ensuremath{\pi}(1 + $y^2$)\}] for y \textgreater{} 0 and = 0 otherwise
\end{enumerate}
\item Find the current through a 0.01 \ensuremath{\muF} capacitor at t = 0 if the voltage across it is (a) 2 sin 2\ensuremath{\pi} $10^6$ t V; (b) - 2 e^-(10)7 t V; (c) - 2 e-(10)7 tsin 2\ensuremath{\pi} $10^6$ t V.
\begin{enumerate}[label=(\alph*)]
    \item 0.1256 A, 0.2 A, 0.1256 A
    \item 0.1256 A, 0.3 A, 0.1256 A
    \item 0.2345 A, 0.1 A, - 0.4567 A
    \item 0.2 A, 0.1 A, 0.2 A
    \item 0 A, 0.1 A, 0 A
\end{enumerate}
\item A certain load to be driven at 1750 r/min requires a torque of 60 lb. ft. What horsepower will be required to drive the load ?
\begin{enumerate}[label=(\alph*)]
    \item 15 hp
    \item 20 hp
    \item 10 hp
    \item 35 hp
    \item 5 hp
\end{enumerate}
\item The moving coil-meters, damping is provided by
\begin{enumerate}[label=(\alph*)]
    \item the resistance of the coil.
    \item the inertia of the coil.
    \item the coil spring attached to the moving.
    \item the electric charge in the coil.
    \item the magnetic field created by the coil.
\end{enumerate}
\item The current through a resistor of 2\ensuremath{\Omega} is given by i(t) = $cos^2$\ensuremath{\pit}. Find the energy dissipated in the resistor from t\_0 = 0 to t\_1 = 5 sec.
\begin{enumerate}[label=(\alph*)]
    \item 5.75 J
    \item 7.50 J
    \item 4.75 J
    \item 1.25 J
    \item 3.75 J
\end{enumerate}
\end{enumerate}

\section*{True / False (10 marks)}
\textit{Answer True or False}
\begin{enumerate}[label=\arabic*.]
\setcounter{enumi}{5}
\item True or False: The radiation resistance of a single-turn circular loop with a circumference of (1/4) wavelength is 0.77 ohm.
\item True or False: One liter-atmosphere of work is equivalent to 106.4 joules.
\item True or False: Newton's Law of cooling for convective heat transfer is represented by the equation q = hAΔT.
\item True or False: A D-flip-flop is transparent when the output is triggered by the rising edge of the clock signal.
\item True or False: The electrostatic energy of a point charge q inside a conducting shell is zero due to shielding effects of the shell.
\end{enumerate}

\section*{Long Form Response (20 marks)}
\textit{Answer each question with a clear and complete explanation.}
\begin{enumerate}[label=\arabic*.]
\setcounter{enumi}{10}
\item Discuss the principles involved in introducing air through a nozzle into water to create bubbles of 2 mm diameter. Calculate the necessary pressure difference between the air at the nozzle and the surrounding water, considering the surface tension of water given as σ = 72.7 \ensuremath{\times} 10^-3 N m^-1.
\item Analyze the function F(s) = (3 s + 1) / \{(s - 1)($s^2$ + 1)\} and determine its inverse Laplace transform. Explain your reasoning and the steps taken in your analysis.
\end{enumerate}

\vfill
\noindent\textit{End of Paper}
\end{document}